\PassOptionsToPackage{AutoFakeBold}{xeCJK}
\documentclass[a4paper,zihao=-4,linespread=1.25]{ctexart}

%===========================================================
% 宏包导入
%===========================================================
\usepackage{geometry}       % 页面设置
\usepackage{fontspec}       % 字体设置
\usepackage{titlesec}       % 标题格式
\usepackage{graphicx}       % 图片支持
\usepackage{booktabs}       % 三线表
\usepackage{caption}        % 图表标题
\usepackage{bicaption}      % 双语标题
\usepackage{float}          % 浮动体控制
\usepackage{amsmath}        % 数学公式
\usepackage{amssymb}        % 数学符号
\usepackage{listings}       % 代码块
\usepackage{xcolor}         % 颜色
\usepackage{fancyhdr}       % 页眉页脚
\usepackage[hidelinks]{hyperref} % 超链接
\usepackage{setspace}       % 行距细调
\usepackage{titletoc}       % 目录格式

%===========================================================
% 页眉页脚与页码
%===========================================================
% 需求：页码位于页面底端居中；页面上方不显示任何内容
\renewcommand{\headrulewidth}{0pt}
\renewcommand{\footrulewidth}{0pt}
\fancypagestyle{plain}{
    \fancyhf{}
    \fancyfoot[C]{\zihao{-4}\thepage}
}

%===========================================================
% 注释（脚注）样式：圆圈加数字右上标
%===========================================================
\renewcommand{\thefootnote}{\textcircled{\scriptsize\arabic{footnote}}}

%===========================================================
% 目录设置 (titletoc)
%===========================================================
% 目录页标题设置
\renewcommand{\contentsname}{目录}

% 重新定义 \tableofcontents 以居中标题
\let\oldtableofcontents\tableofcontents
\renewcommand{\tableofcontents}{%
    \begingroup
    % 临时修改 section 格式以使目录标题居中
    \ctexset{section/format={\centering\zihao{-3}\bfseries\songti}}
    \oldtableofcontents
    \endgroup
}

% 目录内容样式设置
% section: 左对齐, 0缩进, 标题后加点, 引导线为点
\titlecontents{section}[0pt]{\songti\zihao{-4}}
    {\thecontentslabel.\quad} % 有序号
    {}                        % 无序号 (如参考文献)
    {\titlerule*[8pt]{.}\contentspage}

% subsection: 左缩进 2em
\titlecontents{subsection}[2em]{\songti\zihao{-4}}
    {\thecontentslabel\quad}
    {}
    {\titlerule*[8pt]{.}\contentspage}

% subsubsection: 左缩进 4em
\titlecontents{subsubsection}[4em]{\songti\zihao{-4}}
    {\thecontentslabel\quad}
    {}
    {\titlerule*[8pt]{.}\contentspage}
%===========================================================
% 参考文献设置
%===========================================================
% 使用 gbt7714 宏包符合国标
% super: 上标引用 [1]
% square: 方括号
% comma: 逗号分隔
% sort&compress: 排序并压缩连续序号
\usepackage[super,square,comma,sort&compress]{natbib}
\bibliographystyle{gbt7714-numerical}

% 兼容性设置：如果 natbib 与 ctex 冲突，可调整引用样式
% 这里使用 natbib 的标准上标设置
\setcitestyle{super,square,comma}

%===========================================================
% 页面设置
%===========================================================
% 上边距：2.54cm；下边距：2.54cm；左边距 3.17cm；右边距：3.17cm
\geometry{top=2.54cm, bottom=2.54cm, left=3.17cm, right=3.17cm}

%===========================================================
% 字体设置
%===========================================================
% 英文字体：Times New Roman
\setmainfont{Times New Roman}

% 中文字体：ctex 会自动调用 windows 字体 (SimSun, SimHei, KaiTi)
% 配合 AutoFakeBold 选项解决粗体问题

%===========================================================
% 标题格式设置 (严格按照规范)
%===========================================================
% 各级标题：宋体小四号加粗打印
% 理科以“1”、“1.1”、“1.1.1”等数字依次标出，左起打印

\ctexset{
    section = {
        format = \zihao{-4}\bfseries\songti, % 宋体小四加粗
        number = \arabic{section},
        name = {,},
        aftername = \hspace{0.5em},
        indent = 0pt, % 左起打印
        beforeskip = 1ex,
        afterskip = 1ex
    },
    subsection = {
        format = \zihao{-4}\bfseries\songti,
        number = \arabic{section}.\arabic{subsection},
        indent = 0pt,
        beforeskip = 1ex,
        afterskip = 1ex
    },
    subsubsection = {
        format = \zihao{-4}\bfseries\songti,
        number = \arabic{section}.\arabic{subsection}.\arabic{subsubsection},
        indent = 0pt,
        beforeskip = 1ex,
        afterskip = 1ex
    }
}

%===========================================================
% 图表标题设置
%===========================================================
% 图题：五号楷体，居中，图下
% 表题：五号宋体加粗，居中，表上
% 英文标题：Times New Roman 五号，居中
% 图号/表号：按全文出现顺序以阿拉伯数字编排 (1, 2, 3...)

\DeclareCaptionFont{kaitifive}{\zihao{5}\kaishu}
\DeclareCaptionFont{songtifivebold}{\zihao{5}\bfseries\songti}
\DeclareCaptionFont{englishfont}{\zihao{5}\rmfamily}

% 基本标题设置
\captionsetup{
    labelsep=quad,
    skip=6pt
}

% 图设置
\captionsetup[figure]{
    font=kaitifive, 
    name={图}
}

% 表设置
\captionsetup[table]{
    font=songtifivebold, 
    name={表}
}

% 双语标题设置 (bicaption)
% 关键修正：不使用 bi-second 选项组，而是直接配置
\captionsetup[figure][bi-second]{name={Fig.}, font=englishfont}
\captionsetup[table][bi-second]{name={Table}, font=englishfont}

% 确保图表编号连续 (不带章节号)
\usepackage{chngcntr}
\counterwithout{figure}{section}
\counterwithout{table}{section}

%===========================================================
% 代码块设置
%===========================================================
\lstset{
    basicstyle=\ttfamily\footnotesize,
    breaklines=true,
    frame=single,
    numbers=left,
    numberstyle=\tiny,
    keywordstyle=\color{blue},
    commentstyle=\color{gray},
    stringstyle=\color{green!50!black}
}

%===========================================================
% 图片路径
%===========================================================
\graphicspath{{../docs/figures/}}

%===========================================================
% 论文题目样式
%===========================================================
\newcommand{\thesistitle}[1]{
    \begin{center}
        {\zihao{-3}\bfseries\songti #1} % 宋体小三号加粗，居中
    \end{center}
}

%===========================================================
% 文档开始
%===========================================================
\begin{document}

% 封面、摘要等前置部分不计入页码或使用罗马数字（根据规范“论文页码从正文序论部分开始”）
\pagenumbering{gobble} % 暂时关闭页码
\pagestyle{empty} % 前置部分不显示页眉页脚

%-----------------------------------------------------------
% 封面 (模拟)
%-----------------------------------------------------------
\begin{titlepage}
    \centering
    \vspace*{2cm}
    {\zihao{1}\bfseries 毕业论文（设计）\par}
    \vspace{3cm}
    
    {\zihao{-3}\bfseries MP-CycleGAN：基于多感知驱动的非配对水下图像增强方法\par}
    
    \vspace{4cm}
    
    {\zihao{-4}\songti
    \begin{tabular}{rl}
        学校： & \underline{\makebox[6cm]{云南师范大学}} \\
        学院： & \underline{\makebox[6cm]{信息学院}} \\
        专业： & \underline{\makebox[6cm]{计算机科学与技术}} \\
        学生姓名： & \underline{\makebox[6cm]{XXX}} \\
        学号： & \underline{\makebox[6cm]{XXXXXXXX}} \\
        指导教师： & \underline{\makebox[6cm]{XXX}} \\
        完成日期： & \underline{\makebox[6cm]{2026年5月}} \\
    \end{tabular}\par}
    
    \vfill
\end{titlepage}

%-----------------------------------------------------------
% 目录 (根据规范未明确列出，但保留以方便查阅，若不需要可注释)
%-----------------------------------------------------------
\newpage
\tableofcontents
\newpage

%-----------------------------------------------------------
% 正文：页码从这里开始
%-----------------------------------------------------------
\newpage
\pagenumbering{arabic} % 从正文开始使用阿拉伯数字
\setcounter{page}{1}   % 重置页码为 1
\pagestyle{plain} % 页码位于页面底端居中，页眉不显示

% 正文文本：宋体小四号，每段首行缩进两格
\zihao{-4}\songti
\setlength{\parindent}{2em}

%-----------------------------------------------------------
% 正文首页：题目 + 摘要 + 关键词（与示例排版一致）
%-----------------------------------------------------------
\thesistitle{MP-CycleGAN：基于多感知驱动的非配对水下图像增强方法}

\vspace{\baselineskip} % 题目下空一行

% 摘要：另起一行，首行缩进两格排印“摘要：”，宋体小四号加粗；其余内容宋体小四号，单倍行距
\noindent\hspace{2em}{\bfseries 摘要：}
{\setstretch{1.0}
水下图像增强是海洋工程、水下机器人与水下视觉任务中的关键预处理步骤。受水体吸收与散射影响，真实水下图像普遍存在蓝绿偏色、对比度降低与雾化等问题，直接削弱下游检测、识别与测量的可靠性。近年来，基于深度学习的监督增强方法在合成或成对数据上取得了显著效果，但其依赖难以获取的像素级配对数据，且在真实跨域场景中容易受到域差异影响而性能下降。

为缓解配对数据依赖并提升跨域鲁棒性，本文在 CycleGAN 的非配对图像到图像转换框架上，提出一种物理与感知双重驱动的无监督水下增强方法 \textbf{MP-CycleGAN}（Multi-Perceptual CycleGAN）。“MP”在本文中特指多感知约束，即同时从颜色感知、结构感知与语义感知三个层面对生成结果进行约束：基于灰世界假设的颜色物理约束用于抑制水下蓝绿偏色，基于梯度一致性的结构约束用于缓解风格迁移导致的结构漂移，基于预训练 VGG-16 特征的感知约束用于保留语义与纹理细节。其核心逻辑不是单纯叠加多个损失项，而是利用不同感知层面的互补信息，共同缩小非配对映射的可行解空间，使增强结果在颜色恢复、结构保真与视觉自然性之间取得更稳定的平衡。本文在统一推理与评测脚本下，对 EUVP 与 UIEB 数据集进行同域评测与跨域泛化测试，并采用 bootstrap 置信区间报告关键指标，提升结论可复核性。

实验结果表明：在 EUVP \texttt{test\_samples}（matched=200）上，CycleGAN（22.82 dB / 0.783）与 MP-CycleGAN（22.91 dB / 0.795）整体表现接近且 MP-CycleGAN 在像素一致性指标上略优；在 UIEB 跨域测试（matched=890）中，MP-CycleGAN 的 UIQM 达到 10.50，与 CycleGAN 的 10.48 基本一致。消融实验进一步验证了灰世界、结构与感知约束对指标的影响趋势，为 MP-CycleGAN 的损失设计提供了定量支撑。
}

\vspace{\baselineskip}

% 关键词：摘要下另起一行，首行缩进两格排印“关键词：”，宋体小四号加粗；词条宋体小四号
\noindent\hspace{2em}{\bfseries 关键词：}水下图像增强；非配对学习；CycleGAN；灰世界；感知损失；结构一致性

\vspace{\baselineskip}

\section{绪论}

\subsection{研究背景与意义}
地球表面约 71\% 被海洋覆盖，海洋蕴含丰富的资源与生态信息。水下图像与视频作为水下机器人导航、海洋资源勘探、生态监测及海底工程检测等任务的核心信息载体，其质量直接决定了后端视觉系统的可靠性与运行效率。然而，在实际作业中，水下视觉系统常受限于低照度、悬浮颗粒、复杂光照与水体类型差异等因素，图像退化现象普遍且程度不一，表现为颜色偏移、对比度下降、细节模糊与雾化等问题。

与空气成像规律不同，光在水介质中的传播同时受到“波长依赖性吸收衰减”与“散射效应”的双重影响：长波长的红光在水中衰减更快，导致图像整体偏蓝或偏绿；后向散射引起的雾化效应又会显著降低全局对比度，掩盖关键边缘与纹理细节。上述退化不仅背离人类主观视觉习惯，更会因特征表达弱化而削弱下游视觉算法（如目标检测、语义分割与三维重建）的鲁棒性。因此，研发高性能的水下图像增强算法以恢复图像色彩与清晰度，不仅具有迫切的工程应用价值，也是计算机视觉与图像处理领域的重要研究课题。

在应用层面，水下图像增强通常处于“视觉链路”的前端，其输出质量会在后续模块中被放大：增强不足会使检测与分割结果不稳定，增强过度则可能引入伪影导致误检或边界错分。因此，增强方法不仅需要在观感上更自然，也需要在结构保真与语义一致性方面具备可靠保障。
\begin{itemize}
    \item \textbf{吸收}：长波（红光）衰减更快，导致整体偏蓝/偏绿；
    \item \textbf{散射}：后向散射引入雾化，降低对比度并破坏细节；
    \item \textbf{光照不均}：光源方向与距离变化导致局部亮度不稳定。
\end{itemize}

水下增强的目标是在尽量保留内容结构的前提下，改善颜色与对比度，使图像更适合视觉感知与后续任务。

\subsection{研究现状}
目前，水下图像增强研究主要形成了传统先验、监督学习与无监督学习三条演进路径。

（1）\textbf{基于物理模型与统计先验的方法}。此类方法依托大气散射模型、光学成像模型或统计规律进行去雾与颜色校正，例如暗通道先验（DCP）及其变体等思路~\cite{5,6}，以及基于多尺度融合与 Retinex 理论的增强框架。它们通常无需训练数据、可解释性较强，但对场景假设依赖明显（例如透射率分布、光照模型或颜色先验），在水体浑浊、光照不均或存在强后向散射时容易出现过增强、局部失真或颜色不自然等问题。

（2）\textbf{基于监督学习的方法}。得益于深度学习的发展，WaterNet~\cite{1}、FUnIE-GAN~\cite{2}、UColor~\cite{3} 等模型通过在配对或合成配对数据上的训练取得了显著进展，能够在特定数据分布下获得较高的客观指标与较好的视觉效果。然而，该方向高度依赖配对数据质量与覆盖范围。由于真实水下环境中获取严格像素对齐的“退化—清晰”图像对极具挑战，且合成数据与真实物理场景间存在分布偏移（Domain Gap），模型在陌生水域与复杂成像条件下往往出现明显性能退化。

（3）\textbf{基于无监督与弱监督的方法}。为摆脱对配对数据的依赖，CycleGAN~\cite{4} 利用循环一致性在非配对样本间建立映射，在工程上更具实用潜力。但原生 CycleGAN 在水下任务中容易产生内容伪影（Artifacts）、颜色过度偏移或结构畸变；同时，全参考指标（如 PSNR/SSIM）与感知类指标（如 UIQM/UCIQE）之间往往存在优化权衡，如何兼顾像素一致性、结构保真与主观观感仍是亟待解决的难点。

\subsection{研究问题与挑战}
综合上述现状，水下图像增强在数据、模型与评价三个层面均存在挑战。
\begin{itemize}
    \item \textbf{退化机理复杂且变化大}：不同水体与成像距离下，吸收与散射程度差异显著，单一先验难以覆盖全部场景，模型需要具备一定自适应能力。
    \item \textbf{成对数据匮乏与域差异}：真实环境中获取严格像素对齐的配对数据极具挑战，合成数据虽可扩充样本但易引入 Domain Gap，导致跨域泛化不稳定。
    \item \textbf{颜色漂移与过增强风险}：无监督风格迁移易产生不自然色调或局部过度白化/饱和，影响观感并引入伪影。
    \item \textbf{结构漂移与伪影问题}：模型可能在纹理与边缘处产生结构畸变、振铃或不真实纹理，影响下游任务可靠性。
    \item \textbf{指标冲突与评价困难}：PSNR/SSIM 强调与参考的一致性，而 UIQM/UCIQE 更偏向感知质量与对比度等属性，二者在跨域条件下常出现不一致，难以用单一指标全面评价。
\end{itemize}

\subsection{本文工作与贡献}
本文主要工作如下：
\begin{enumerate}
    \item 在 CycleGAN 基础上提出 MP-CycleGAN，并将“MP”明确为颜色感知、结构感知与语义感知三类互补约束的统一建模，引入灰世界物理约束、结构梯度一致性与深层感知约束，缓解偏色与结构漂移。
    \item 在统一评测脚本下对 EUVP 与 UIEB 进行同域/跨域评测，报告 PSNR/SSIM 的 bootstrap 置信区间以增强可复核性。
    \item 通过小预算消融实验分析不同自监督项的作用，给出定量趋势支撑。
\end{enumerate}

\subsection{论文结构}
本文其余章节安排如下：第 2 章介绍水下图像增强的相关研究进展，并梳理 CycleGAN 的基本原理、适用优势与在水下场景中的局限性，为本文方法设计提供理论基础；第 3 章详细阐述 MP-CycleGAN 的整体框架、多感知自监督损失构成及其协同机制，说明本文方法如何在非配对设置下兼顾颜色恢复、结构保持与视觉自然性；第 4 章介绍实验数据集、评测协议、评价指标与统计设置，并通过同域对比、跨域测试、消融实验和可视化结果对方法有效性进行验证与分析；第 5 章总结全文工作，概括本文方法的主要结论、局限性及后续改进方向。

\section{相关工作与技术基础}

本章从“先验方法—监督方法—非监督方法”的发展脉络出发，简要回顾水下图像增强的代表性思路，并重点介绍 CycleGAN 的关键机制与局限性，为后续提出的 MP-CycleGAN 损失设计与实验对比提供背景与对照基线。

\subsection{传统水下增强方法}
传统方法多依赖先验与物理模型，例如暗通道先验（DCP）及其水下变体（UDCP）、Retinex、融合增强等。其优点是无需训练数据、可解释性强，但在复杂水体、光照变化与散射程度差异显著时，容易出现过增强或颜色失真。

\subsection{监督学习方法}
监督方法通过成对数据学习从退化到清晰的映射，代表性方法包括 WaterNet~\cite{1}、UColor~\cite{3}、FUnIE-GAN~\cite{2} 等。该类方法通常能取得较高的客观指标，但对配对数据强依赖，且跨域泛化能力受限。

\subsection{无监督方法与 CycleGAN}
CycleGAN~\cite{4} 通过对抗学习与循环一致性在非配对数据上实现域间转换，显著降低了配对数据需求。但标准 CycleGAN 在水下增强中仍可能产生颜色漂移、纹理伪影或结构被破坏。基于上述分析，下一章将在 CycleGAN 框架之上引入物理与感知层面的约束，给出更稳定、可控且适用于水下增强任务的 MP-CycleGAN 设计。

\section{方法：MP-CycleGAN}

本章详细介绍 MP-CycleGAN 的整体框架与训练目标函数。首先给出生成器与判别器组成及推理流程；随后在 CycleGAN 基础损失之上引入物理与感知双重驱动的自监督约束，并讨论各损失项的协同关系与权重设置策略，解释本文方法为何能够在保持非配对学习优势的同时提升增强的稳定性与可控性。

\subsection{整体框架}
MP-CycleGAN 采用双向循环生成架构，包含两个生成器 $G: A \to B$（水下 $\to$ 清晰）、$F: B \to A$（清晰 $\to$ 水下），以及两个判别器 $D_A, D_B$。其中 $A$ 为水下域，$B$ 为清晰域。训练时通过对抗损失让生成结果逼近目标域分布，通过循环一致性约束保证内容可逆，从而在非配对条件下学习映射。

与监督方法不同，非配对设置下无法对同一场景的水下退化图像与清晰参考图像做像素级对齐，因此模型必须同时解决两类难题：一是学习域间风格迁移（颜色分布、对比度与雾化程度），二是尽可能保留输入内容结构（几何形状、边缘与语义布局）。CycleGAN 的核心思想是用“对抗学习”实现风格匹配、用“循环一致性”约束内容保真；本文提出的 MP-CycleGAN 则进一步从物理先验与感知特征角度增强这一约束，使生成器在提升观感的同时更可控、稳定。

在推理阶段，系统仅使用 $G: A\to B$ 将真实水下图像映射到增强域输出结果；$F$ 与判别器仅用于训练阶段构建约束与对抗博弈，因此本文方法在部署时与常规 CycleGAN 的推理路径一致，工程集成难度较低。

\begin{figure}[H]
    \centering
    \includegraphics[width=0.9\textwidth]{mp_cyclegan_architecture.png}
    \bicaption{MP-CycleGAN 整体架构示意图}{Schematic diagram of MP-CycleGAN architecture}
    \label{fig:arch}
\end{figure}

\subsection{基础损失（CycleGAN）}
基础框架沿用 CycleGAN 的对抗学习与循环一致性约束。设 $x\sim p_A(x)$ 为水下域样本、$y\sim p_B(y)$ 为清晰域样本：

从优化角度看，CycleGAN 是一个双向的最小最大问题：生成器试图“欺骗”判别器以匹配目标域分布，判别器试图区分真实样本与生成样本。由于无配对监督信号，若仅使用对抗损失，生成器可能出现模式崩塌或在内容上产生不可逆变形；因此必须结合循环一致性与（可选的）恒等映射约束来稳定训练并保留语义结构。

1) \textbf{对抗损失}：约束生成分布逼近目标域分布。本文默认采用最常用的 LSGAN 形式（实现细节与原仓库保持一致）：
\begin{equation}
\mathcal{L}_{GAN}(G,D_B)=\mathbb{E}_{y}[(D_B(y)-1)^2]+\mathbb{E}_{x}[D_B(G(x))^2],
\end{equation}
并对 $F,D_A$ 亦定义同类损失。

其中 $D_B(\cdot)$ 表示判别器对“属于清晰域”的置信输出。LSGAN 将原始 GAN 的对数似然形式替换为最小二乘目标，通常能在训练早期提供更稳定的梯度，减轻判别器过强导致的梯度消失问题，从而提升非配对转换训练的稳定性。由于本文方法在 CycleGAN 框架上扩展损失项，采用与基线一致的对抗形式也便于进行公平对比与消融分析。

2) \textbf{循环一致性损失}：约束映射近似可逆，减少内容漂移：
\begin{equation}
\mathcal{L}_{cyc}(G,F)=\mathbb{E}_{x}[\|F(G(x))-x\|_1]+\mathbb{E}_{y}[\|G(F(y))-y\|_1].
\end{equation}

循环一致性从“可逆性”角度约束内容保真：若 $G$ 将水下图像映射到清晰域后再经 $F$ 映射回水下域，应尽量重构原图像。该约束相当于在无配对条件下提供一种自监督信号，能显著抑制几何结构被破坏、目标消失或出现不相关纹理等问题。本文采用 $L_1$ 距离作为重构误差，是因为 $L_1$ 相较 $L_2$ 更不易产生过度平滑，有助于保留边缘与细节。

3) \textbf{恒等损失（可选）}：在部分设置中用于稳定颜色映射（本文主结果未依赖该项提升）：
\begin{equation}
\mathcal{L}_{idt}(G,F)=\mathbb{E}_{y}[\|G(y)-y\|_1]+\mathbb{E}_{x}[\|F(x)-x\|_1].
\end{equation}

恒等损失的直觉是：当输入已经来自目标域时，生成器不应对其做不必要的风格修改。例如，将清晰域图像输入 $G$ 时应尽量保持不变，从而抑制颜色空间发生非预期偏移。对于水下增强而言，恒等损失往往能缓解过度校色与色调漂移，但也可能限制增强强度、导致对比度提升不足，因此本文将其作为可选项并在实验中未将其作为主要改进来源。

\subsection{多感知自监督损失（MP）}
在基础损失之上，本文引入物理与感知层面的自监督约束，形成总目标：
\begin{equation}
\mathcal{L}_{total}=\mathcal{L}_{GAN}+\lambda_{cyc}\mathcal{L}_{cyc}+\lambda_{idt}\mathcal{L}_{idt}+\mathcal{L}_{MP}
\end{equation}
其中
\begin{equation}
\mathcal{L}_{MP}=\lambda_{gray}\mathcal{L}_{gray}+\lambda_{struct}\mathcal{L}_{struct}+\lambda_{perc}\mathcal{L}_{perc}+\lambda_{color}\mathcal{L}_{color}.
\end{equation}

这里“MP”并非泛指多个损失项的简单堆叠，而是指面向增强结果的多层次感知建模：颜色感知关注通道平衡与色偏校正，结构感知关注边缘、纹理与几何轮廓保持，语义感知关注高层内容一致性与自然视觉风格。多感知损失的设计目标是为非配对学习提供“更接近水下成像机理与人类视觉感知”的额外约束。仅依赖对抗与循环一致性时，模型容易在颜色、对比度与纹理上出现多解；加入物理先验与感知特征约束，可以缩小可行解空间，使模型更倾向于输出自然、结构稳定且语义一致的增强结果。与此同时，各项损失的权重需要进行平衡：过强的颜色约束可能导致过白或过饱和，过强的结构约束可能限制增强幅度，过强的感知约束则可能压制颜色校正效果。因此本文采用可配置的权重形式，并在后续消融实验中分析不同组合的影响趋势。

\subsubsection{灰世界损失（物理约束）}
基于灰世界假设（平均反射率为灰色），约束生成图像 RGB 通道均值趋于平衡：
\begin{equation}
\mathcal{L}_{gray}(\hat{I})=\frac{1}{3}\sum_{c\in\{r,g,b\}}\left|\mu_c(\hat{I})-\mu_{gray}(\hat{I})\right|,\quad
\mu_{gray}(\hat{I})=\frac{\mu_r(\hat{I})+\mu_g(\hat{I})+\mu_b(\hat{I})}{3}.
\end{equation}
该项主要用于抑制水下蓝/绿偏色。

水下环境中红光衰减更快，导致图像整体偏蓝或偏绿。灰世界假设认为自然场景在平均意义下应接近“灰色”，即三通道统计不应长期失衡。尽管该假设并非对所有场景严格成立，但它为颜色校正提供了简单有效的先验约束，能够在无配对监督下引导模型避免产生不自然的色调漂移。与直接进行白平衡不同，本文将该先验以可微损失形式嵌入训练，使生成器在学习目标域风格的同时受到物理合理性约束。

\subsubsection{结构一致性损失（梯度约束）}
为缓解风格迁移带来的结构漂移，采用基于一阶差分的梯度一致性约束：
\begin{equation}
\mathcal{L}_{struct}(I,\hat{I})=\left\|\nabla_x \hat{I}-\nabla_x I\right\|_1+\left\|\nabla_y \hat{I}-\nabla_y I\right\|_1.
\end{equation}
该项直接约束边缘与纹理变化，有助于结构保真。

梯度算子 $\nabla_x,\nabla_y$ 可通过相邻像素差分实现，能够突出边缘与细节纹理。对水下增强任务而言，“风格迁移”主要应该改变颜色与对比度，而不应改变几何结构；因此结构一致性损失在特征层面约束“变化发生在合理方向”，对抑制边缘错位、纹理拉伸与局部结构扭曲具有直接效果。该损失与循环一致性互补：循环一致性约束整体可逆性，结构一致性更关注局部边缘与纹理稳定，二者共同提升结构保真。

\subsubsection{感知损失（深层特征约束）}
利用 ImageNet 预训练的 VGG-16 特征提取器 $\phi$，在特征空间约束输入与生成结果的语义一致性：
\begin{equation}
\mathcal{L}_{perc}(I,\hat{I})=\left\|\phi(I)-\phi(\hat{I})\right\|_1.
\end{equation}
本文复现实验中采用 \texttt{perceptual\_layer=16}，以平衡语义与细节。

感知损失通过预训练网络的中间特征表达捕捉更高层次的语义与纹理统计，相比像素空间损失更能容忍光照与颜色变化而强调“内容一致”。在非配对增强中，对抗损失可能驱动模型生成“看起来像目标域”的纹理，但其中部分纹理可能并不属于原场景内容，从而引入伪影；感知约束能够在一定程度上抑制这种非自然纹理，提升结果的真实性与语义一致性。感知层选择存在权衡：较浅层更关注边缘与纹理，较深层更关注语义结构；本文采用的 \texttt{perceptual\_layer=16} 处于中等深度，以兼顾细节与语义稳定。

\subsubsection{颜色统计一致性损失（全局统计约束）}
为进一步约束全局颜色分布，加入通道均值与标准差匹配项：
\begin{equation}
\mathcal{L}_{color}(I,\hat{I})=\left\|\mu(\hat{I})-\mu(I)\right\|_1+\left\|\sigma(\hat{I})-\sigma(I)\right\|_1.
\end{equation}
该项在消融配置 E 中启用，用于观察对无参考指标的影响趋势。

颜色统计一致性从“全局一阶与二阶统计”角度约束生成图像与输入图像的颜色变化幅度。其作用并非简单保持输入颜色，而是防止生成器在对抗训练中产生过度的全局偏移（例如整体发白或过饱和），从而提升增强结果的稳定性与可控性。在实践中，该项往往对 UIQM 等与对比度/饱和度相关的无参考指标更敏感，但也可能牺牲部分与参考一致性相关的 PSNR/SSIM，因此本文将其作为可选项并通过消融实验展示其影响趋势。

\subsubsection{损失项协同关系与权重设置}
从优化目标角度看，各损失项具有互补性：灰世界损失主要从物理先验出发约束颜色通道统计，适合缓解蓝绿偏色；结构一致性损失直接约束梯度变化，能在风格迁移过程中保留边缘与纹理，减少结构漂移；感知损失通过深层语义特征对齐，抑制对抗训练中可能出现的非自然纹理与伪影，使增强结果在保持语义一致性的同时更符合视觉感知；颜色统计一致性损失进一步约束全局分布，有助于改善对比度与色彩饱和度相关的无参考指标，但也可能带来与参考像素不一致的风险。

在实现上，各项权重通过命令行参数独立控制，便于在不同数据分布与应用场景下进行调参。本文建议先在同域 matched 条件下以 PSNR/SSIM 与视觉效果为主导确定结构与感知项的合理范围，再逐步加入灰世界与颜色统计项以改善色偏与整体观感，避免一次性引入过多强约束导致训练不稳定或出现过增强现象。

\subsection{工程实现要点}
本项目在 CycleGAN 模型中实现上述自监督项，并通过命令行参数控制权重：
\begin{itemize}
    \item \texttt{--lambda\_gray}：灰世界损失权重
    \item \texttt{--lambda\_struct}：结构一致性损失权重
    \item \texttt{--lambda\_perceptual}：感知损失权重
    \item \texttt{--lambda\_color}：颜色统计一致性损失权重
    \item \texttt{--perceptual\_layer} / \texttt{--perceptual\_weights}：控制 VGG-16 特征层深度与权重来源
\end{itemize}
对应实现位于 \texttt{models/cycle\_gan\_model.py}。

基于上述方法设计，下一章将给出统一的评测协议与指标体系，并通过同域评测、跨域泛化测试与消融实验验证 MP-CycleGAN 的有效性与适用边界。

\section{实验设计与结果分析}

本章围绕“是否有效提升水下图像质量、是否具备跨域泛化能力、各损失项的边际贡献是什么”三个问题展开。具体而言，本文首先给出数据集与评测协议，并说明全参考与无参考指标的含义与适用范围；随后报告同域定量结果与跨域测试结果，辅以可视化对比与颜色分布分析；最后通过消融实验与讨论分析指标权衡、失败模式与工程开销，为结论的可信性与可复核性提供支撑。

\subsection{数据集与评测协议}
\subsubsection{EUVP（同域评测）}
训练：EUVP 的 Unpaired 分支~\cite{2}。
评测：EUVP \texttt{test\_samples}（matched=200），使用输入与 GT 对齐计算 PSNR/SSIM，并统计无参考指标 UCIQE/UIQM。

为保证对比公平性，本文在同一推理与评测流程下生成增强结果：对每张输入水下图像仅使用 $G: A\to B$ 方向输出增强图像，并与 \texttt{test\_samples} 中提供的对齐参考图像逐样本计算全参考指标。由于 matched 子集包含严格对齐的参考图像，它更能反映模型在像素一致性与结构保真方面的表现；同时，无参考指标用于刻画整体视觉质量趋势，二者联合报告更贴近真实应用中的多目标需求。

\subsubsection{UIEB（跨域泛化测试）}
训练：仅在 EUVP 上训练，不使用 UIEB 微调。
评测：UIEB（matched=890）~\cite{1}，用 \texttt{reference-890} 作为参考计算 PSNR/SSIM，同时统计 UCIQE/UIQM。

跨域测试旨在检验模型在未见水体类型与成像条件下的稳健性。UIEB 数据集覆盖了更复杂的色偏、散射与光照分布，其 matched 子集提供参考图像用于计算 PSNR/SSIM。需要强调的是，不同数据集的参考图像获取方式与图像风格可能不同，因此跨域评测中 PSNR/SSIM 与无参考指标之间更容易出现差异，这也为后续分析“指标权衡”提供了客观依据。

\subsection{评价指标}
本文采用两类指标：
\begin{itemize}
    \item \textbf{全参考指标}：PSNR、SSIM~\cite{8}
    \item \textbf{无参考指标}：UCIQE~\cite{11}、UIQM~\cite{12}
\end{itemize}
其中 UIQM 由颜色、清晰度与对比度等分量加权组合得到，不同指标关注点不同，因此可能出现“PSNR/SSIM 提升但 UIQM 不一致”的现象。

在解释指标时，本文采用如下原则：PSNR 对像素级误差敏感，适合衡量增强结果与参考图像在灰度/颜色上的一致性；SSIM 更强调局部结构相似性与对比度一致性，更接近结构保真。UCIQE 与 UIQM 属于无参考质量指标，其优势是可用于缺少参考图像的真实场景评估，但也可能偏好更强的对比度或饱和度提升。因此，本文重点关注同域 matched 条件下 PSNR/SSIM 的变化，并结合无参考指标与可视化结果，综合评估增强效果与潜在副作用（如过增强、颜色漂移与纹理伪影）。

\subsection{统计设置与置信区间}
PSNR/SSIM 的 95\% 置信区间通过 bootstrap 估计（\texttt{bootstrap\_iters=2000}，\texttt{seed=123}）。UCIQE/UIQM 报告为全体样本均值。

具体地，对 matched 测试集中 $N$ 对（输入，参考）样本进行有放回重采样，每次采样得到 $N$ 对样本并计算对应指标均值，重复 \texttt{bootstrap\_iters} 次后取 2.5\% 与 97.5\% 分位作为 95\% 置信区间。该过程可缓解小样本波动对结论的影响，使“改进幅度是否稳定”更具可复核性；相关脚本与输出文件路径已在附录给出。

\subsection{主结果（EUVP 同域评测）}

\begin{table}[H]
    \centering
    \bicaption{EUVP test\_samples 同域评测结果（matched=200）}{Quantitative results on EUVP test\_samples (matched=200)}
    \label{tab:euvp_main}
    \begin{tabular}{lcccccc}
        \toprule
        方法 & PSNR (dB) & SSIM & UCIQE & UIQM \\
        \midrule
        Identity & 20.33 & 0.794 & 24.88 & 5.77 \\
        Gray-World & 19.42 & 0.784 & 24.27 & 6.72 \\
        CycleGAN & 22.82 & 0.783 & 25.93 & 5.33 \\
        MP-CycleGAN & \textbf{22.91} & \textbf{0.795} & 25.82 & 5.41 \\
        \bottomrule
    \end{tabular}
    \vspace{0.5em}
    \par{\zihao{5}\songti 注：PSNR/SSIM 为均值，置信区间略。}
\end{table}

表 \ref{tab:euvp_main} 给出了同域 matched=200 条件下的定量对比。与 CycleGAN 相比，MP-CycleGAN 在 PSNR 上提升 0.09 dB（22.82 $\to$ 22.91），在 SSIM 上提升 0.012（0.783 $\to$ 0.795），体现出引入物理与感知约束后对像素一致性与结构保真的正向作用。另一方面，Gray-World 作为纯物理先验基线，其 UIQM 明显更高（6.72），但 PSNR/SSIM 低于 Identity 与学习型方法，说明单纯的颜色均衡可能改善主观观感，却不足以恢复与参考图像一致的细节与对比度。该结果强调了“物理先验与数据驱动学习互补”的必要性：物理先验能约束色偏与全局统计，而对抗学习与循环一致性更擅长恢复纹理与局部结构。

需要说明的是，本文同域指标提升幅度相对有限，主要有三方面原因。其一，CycleGAN 基线本身已具备较强的非配对映射能力，在 matched 子集上的性能处于较高水平，因而后续改进更多体现为“细化校正”而非大幅跃升；其二，本文新增约束主要作用于颜色稳定性、结构保真与伪影抑制，这些收益往往首先体现在边缘过渡更平滑、色偏更轻和视觉自然性更好，而不一定被 PSNR 这类像素误差指标充分放大；其三，matched 参考图像与模型输出在亮度风格、颜色分布上仍可能存在系统性差异，即使增强结果主观上更自然，也未必转化为显著的全参考增益。因此，本文更强调“小幅但稳定”的量化改进，以及定性结果中更清晰的结构边界、更自然的白平衡和更少的局部伪影。

\subsection{消融实验（小预算跨域设置）}
消融实验用于分析损失项叠加的指标趋势。考虑到完整训练与评测成本较高，本文在跨域评测设置下选取较小的实验预算快速对比不同损失组合，因此更关注“趋势与相对增益”，而非追求所有指标在单一配置下的绝对最优。

\begin{table}[H]
    \centering
    \bicaption{不同损失组合的性能对比（小预算跨域消融）}{Performance comparison of different loss combinations (ablation study)}
    \label{tab:ablation}
    \begin{tabular}{clcccc}
        \toprule
        ID & 描述 & PSNR & SSIM & UCIQE & UIQM \\
        \midrule
        A & CycleGAN (Baseline) & 16.85 & 0.56 & 27.89 & 4.13 \\
        B & +Gray & 17.04 & 0.56 & \textbf{28.66} & 4.16 \\
        C & +Gray+Struct & 16.67 & 0.65 & 26.85 & 4.18 \\
        D & +Gray+Struct+Perc & \textbf{17.80} & \textbf{0.65} & 28.51 & 4.56 \\
        E & +Color & 16.72 & 0.64 & 27.87 & \textbf{4.76} \\
        \bottomrule
    \end{tabular}
\end{table}

从表 \ref{tab:ablation} 可以观察到不同自监督项对指标侧重点的影响规律：加入灰世界约束（A$\to$B）显著提升 UCIQE（27.89 $\to$ 28.66），表明该物理先验能够有效缓解色偏并提升颜色相关质量；在此基础上加入结构梯度一致性（B$\to$C）使 SSIM 显著上升（0.56 $\to$ 0.65），反映结构约束能在跨域更复杂条件下更稳定地保留边缘与纹理，但其对 UCIQE 的影响并非单调，可能与“更强的结构保真”限制了对比度增强空间有关。进一步加入感知损失（C$\to$D）带来 PSNR 与 UIQM 的同步提升（PSNR 16.67 $\to$ 17.80，UIQM 4.18 $\to$ 4.56），说明深层特征约束能够在不破坏结构的前提下，提高整体感知质量并降低伪影。最后加入颜色统计一致性（D$\to$E）将 UIQM 推至最高（4.76），但 PSNR/SSIM 略有回落，这与无参考指标对全局对比度/饱和度更敏感、而全参考指标更强调与参考像素一致性的差异是一致的。

\begin{figure}[H]
    \centering
    \includegraphics[width=0.9\textwidth]{ablation_study_grid.png}
    \bicaption{消融实验量化对比}{Quantitative comparison of ablation study}
    \label{fig:ablation}
\end{figure}

\subsection{视觉对比与颜色分布分析}

仅依赖单一客观指标可能掩盖增强算法的实际视觉表现，因此本文进一步结合定性结果分析，以从主观感知层面验证增强结果的有效性。首先展示同域条件下的典型视觉样例，用于比较输入图像、基线方法与本文方法在颜色恢复、局部对比度与结构保持方面的差异；随后再从统计分布角度分析增强结果与参考域之间的颜色关系，以增强对实验现象的解释性。

\begin{figure}[H]
    \centering
    \includegraphics[width=0.9\textwidth]{visual_comparison_final_refined.png}
    \bicaption{EUVP 同域视觉对比。（a）Input，（b）CycleGAN，（c）MP-CycleGAN，（d）GT}{Visual comparison on EUVP. (a) Input, (b) CycleGAN, (c) MP-CycleGAN, (d) GT}
    \label{fig:visual}
\end{figure}

由图 \ref{fig:visual} 可以看出，相较于 CycleGAN，MP-CycleGAN 的颜色分布更接近参考图像的自然状态，水下场景中常见的蓝绿偏色得到明显缓解，且前景目标与背景区域之间的亮度层次更加清晰。同时，本文方法在边缘轮廓、细小纹理与局部结构上表现出更好的稳定性，说明结构与感知约束能够在增强色彩的同时抑制风格迁移引入的结构漂移。总体而言，该图表明本文方法并非仅通过提升亮度获得主观改善，而是在颜色校正、对比度恢复与结构保持之间实现了更平衡的增强效果。

\begin{figure}[H]
    \centering
    \includegraphics[width=0.9\textwidth]{color_distribution_analysis.png}
    \bicaption{颜色分布分析}{Color distribution analysis}
    \label{fig:color}
\end{figure}

图 \ref{fig:color} 进一步从统计分布角度验证了上述视觉现象。可以观察到，MP-CycleGAN 在各颜色通道上的均值位置与分布形态整体上更接近参考域，说明模型能够有效减轻由水体吸收与散射造成的颜色偏移；同时，通道分布范围的适度扩展反映出对比度和动态范围的提升。需要指出的是，若分布拉伸过强，则可能带来过饱和或局部细节损失，因此本文在训练中将颜色统计一致性作为可调约束，并结合消融实验对这一权衡进行了讨论。综合图 \ref{fig:visual} 与图 \ref{fig:color} 可知，本文方法在主观视觉质量与颜色统计一致性两个层面均表现出较好的稳定性。

\subsection{更多定性结果展示}
为更全面地覆盖不同水体类型、光照条件与场景内容，本节进一步采用横向拼图方式展示多组增强前后对比。相较于传统逐张排布，横向布局更适合在论文版面中连续观察“输入$\to$增强$\to$参考”或“输入$\to$基线$\to$本文方法”的变化趋势，有利于从颜色恢复、局部对比度提升与结构保持三个层面进行综合判断。下文依次从同域三联对比、基线方法对比与跨域视觉效果三个角度展开分析，以形成与量化实验相互补充的定性证据。

\begin{figure}[htbp]
    \centering
    \includegraphics[width=0.98\textwidth]{../pytorch-CycleGAN-and-pix2pix-master/pytorch-CycleGAN-and-pix2pix-master/results/paper_visuals/euvp_triplets_landscape_6cases.png}
    \bicaption{EUVP 同域测试中的横向三联对比。每组样例从左到右依次为输入图像、MP-CycleGAN 增强结果与参考图像，展示了颜色校正、对比度恢复与结构保持效果。}{Landscape triplet comparison on EUVP. For each case, the images from left to right are the input, MP-CycleGAN result and reference, showing color correction, contrast recovery and structure preservation.}
    \label{fig:qual_triplets}
\end{figure}

从图 \ref{fig:qual_triplets} 可以看出，MP-CycleGAN 在多数样例中能够较明显地削弱蓝绿偏色，提高前景目标与背景区域之间的亮度层次，并使船体边缘、珊瑚轮廓及岩石纹理等细节更加可辨。与参考图像相比，增强结果虽然在整体亮度风格上仍存在轻微差异，但主结构位置、颜色过渡与场景语义保持较好一致，说明本文引入的灰世界先验与结构一致性约束在抑制无监督风格迁移引起的内容漂移方面是有效的。由此可见，本文方法在同域场景下具备较好的颜色恢复能力与结构保真能力。

\begin{figure}[htbp]
    \centering
    \includegraphics[width=0.98\textwidth]{../pytorch-CycleGAN-and-pix2pix-master/pytorch-CycleGAN-and-pix2pix-master/results/paper_visuals/euvp_baseline_vs_mpcgan_landscape_6cases.png}
    \bicaption{EUVP 同域测试中的横向方法对比。每组样例从左到右依次为输入图像、CycleGAN 基线结果与 MP-CycleGAN 结果，用于直观比较本文方法在颜色恢复与细节增强方面的改进。}{Landscape method comparison on EUVP. For each case, the images from left to right are the input, CycleGAN baseline result and MP-CycleGAN result, highlighting the improvements in color restoration and detail enhancement.}
    \label{fig:qual_compare}
\end{figure}

图 \ref{fig:qual_compare} 进一步体现了本文方法相对基线 CycleGAN 的优势。基线方法虽然能够一定程度提升亮度与饱和度，但在部分样例中仍保留明显的蓝绿偏色，或出现局部区域发灰、层次不足的问题。相比之下，MP-CycleGAN 的增强结果通常具有更自然的白平衡、更清晰的边缘过渡以及更稳定的局部对比度，这与前文 PSNR/SSIM 和 UIQM 的提升趋势相互印证。需要说明的是，在极端散射或光照不均场景中，本文方法仍可能出现局部过增强、亮部偏白或细小纹理不够稳定等现象，这也提示后续可进一步引入自适应强度控制机制。总体而言，该图从视觉层面验证了本文改进相较基线模型具有稳定且可感知的增益。

\begin{figure}[htbp]
    \centering
    \includegraphics[width=0.98\textwidth]{../pytorch-CycleGAN-and-pix2pix-master/pytorch-CycleGAN-and-pix2pix-master/results/paper_visuals/uieb_pairs_landscape_6cases.png}
    \bicaption{UIEB 跨域测试中的横向输入/增强结果对照。每组样例从左到右依次为输入图像与 MP-CycleGAN 增强结果，用于展示模型在未配对跨域场景中的主观视觉改善。}{Landscape input-output comparison on UIEB cross-domain testing. For each case, the images from left to right are the input and MP-CycleGAN result, showing the subjective visual improvement under unseen cross-domain conditions.}
    \label{fig:qual_pairs}
\end{figure}

从图 \ref{fig:qual_pairs} 可以观察到，即使在训练域之外的 UIEB 数据上，MP-CycleGAN 仍能够在多数组样例中提升画面可见度，改善暗部细节、缓解绿色或蓝色偏移，并增强主体轮廓的清晰度。这说明本文方法学习到的映射关系并非仅适用于特定训练样本，而具有一定跨域泛化能力。当然，在部分极端场景中仍可观察到增强强度偏大或局部纹理不稳定的现象，这与跨域量化实验中全参考指标和无参考指标存在分歧的结论是一致的。综合来看，横向定性结果为本文关于“同域稳定增益、跨域保持竞争性表现”的结论提供了直观支撑。

\begin{figure}[H]
    \centering
    \includegraphics[width=0.98\textwidth]{../pytorch-CycleGAN-and-pix2pix-master/pytorch-CycleGAN-and-pix2pix-master/results/paper_figures/fig_qualitative_euvp.pdf}
    \bicaption{EUVP 测试集的批量定性对比（论文级整图）}{Large-scale qualitative comparisons on EUVP (paper-style figure)}
    \label{fig:qual_euvp_pdf}
\end{figure}

\subsection{与现有方法对比（参考展示）}
为便于读者直观对比，选取 WaterNet~\cite{1}、UColor~\cite{3} 与 FUnIE-GAN~\cite{2} 的公开结果进行参考展示。

需要说明的是，WaterNet、UColor 与 FUnIE-GAN 等方法多以配对或合成配对数据进行监督训练，在训练数据、损失设计与推理策略上与非配对方法存在差异，因此该对比更多用于提供“量级参考”。在非配对设置下，MP-CycleGAN 在 EUVP matched 测试上的 PSNR/SSIM 表现具有竞争力（表 \ref{tab:sota}），说明引入物理与感知双重约束后，非配对框架也能在像素一致性与结构保真方面达到较高水平。

\begin{table}[H]
    \centering
    \bicaption{EUVP 数据集上的参考对比（仅作展示）\cite{1}}{Reference comparison on EUVP dataset (for display only) \cite{1}}
    \label{tab:sota}
    \begin{tabular}{lccc}
        \toprule
        方法 & 类别 & PSNR (dB) & SSIM \\
        \midrule
        WaterNet & Supervised & 19.81 & 0.86 \\
        UColor & Supervised & 21.86 & 0.89 \\
        FUnIE-GAN & Supervised & 23.50 & 0.92 \\
        CycleGAN（复现） & Unsupervised & 22.82 & 0.783 \\
        MP-CycleGAN（复现） & Unsupervised & \textbf{22.91} & \textbf{0.795} \\
        \bottomrule
    \end{tabular}
\end{table}

\begin{figure}[H]
    \centering
    \includegraphics[width=0.9\textwidth]{sota_comparison_refined.png}
    \bicaption{SOTA 方法性能对比（参考展示）}{Performance comparison of SOTA methods (reference display)}
    \label{fig:sota}
\end{figure}

\subsection{\texorpdfstring{跨域泛化性测试（EUVP $\to$ UIEB）}{跨域泛化性测试（EUVP->UIEB）}}

\begin{table}[H]
    \centering
    \bicaption{UIEB 数据集跨域测试结果（matched=890）}{Cross-domain testing results on UIEB dataset (matched=890)}
    \label{tab:uieb}
    \begin{tabular}{lccc}
        \toprule
        指标 & Input & CycleGAN & MP-CycleGAN \\
        \midrule
        PSNR (dB) & - & \textbf{17.84} & 17.54 \\
        SSIM & - & \textbf{0.626} & 0.624 \\
        UCIQE & 21.70 & \textbf{23.25} & 22.46 \\
        UIQM & 6.80 & 10.48 & \textbf{10.50} \\
        \bottomrule
    \end{tabular}
\end{table}

跨域结果如表 \ref{tab:uieb} 所示。为避免仅凭单一指标得出片面结论，本文同时结合跨域可视化结果进行分析，以考察模型在未知水体条件下的主观增强效果与实际可用性。CycleGAN 在 PSNR/SSIM 与 UCIQE 上略占优势，而 MP-CycleGAN 在 UIQM 上略优。该现象可从“优化目标差异”理解：MP-CycleGAN 的多感知约束倾向于提升对比度、颜色均衡与感知锐度，这更容易被 UIQM 捕捉；但当参考图像的风格与亮度分布与增强结果存在系统性差异时，PSNR/SSIM 可能下降。换言之，跨域场景中“更符合主观感知”的增强结果不一定与参考图像像素分布一致，这也解释了无参考指标与全参考指标在跨域条件下更容易出现不一致。

从鲁棒性角度看，本文方法的优势不在于跨域条件下对所有指标都取得绝对领先，而在于面对未见水体类型、散射强度和光照变化时，仍能较稳定地保持颜色校正、局部对比度提升与主体边界清晰度。也就是说，MP-CycleGAN 更偏向输出“不过度失真且可用于观察”的增强结果，而不是单纯追求与参考图像风格的数值贴合。这种特性对于真实水下应用更具意义，因为实际工程环境往往缺少严格参考图像，算法的稳定性、可解释性与视觉可用性通常比单一指标的领先更重要。

\begin{figure}[htbp]
    \centering
    \includegraphics[width=0.98\textwidth]{../pytorch-CycleGAN-and-pix2pix-master/pytorch-CycleGAN-and-pix2pix-master/results/paper_visuals/uieb_pairs_landscape_6cases.png}
    \bicaption{UIEB 跨域测试的横向定性结果。每组样例从左到右依次为输入图像与 MP-CycleGAN 增强结果，反映了模型在未知水体条件下的泛化能力。}{Landscape qualitative results on UIEB cross-domain testing. For each case, the images from left to right are the input and MP-CycleGAN result, reflecting the generalization ability under unseen water conditions.}
    \label{fig:uieb}
\end{figure}

结合图 \ref{fig:uieb} 的横向定性对比可以进一步看到，MP-CycleGAN 在多数组样例中能够有效提升画面可见度，使原本偏暗、偏绿或层次模糊的区域变得更清晰；尤其在主体边缘、纹理轮廓与局部亮度层次方面，其主观视觉质量改善较为明显。但在极端浑浊或强光照不均情况下，仍可能出现局部过增强、亮部偏白或纹理伪影，提示未来可引入更强的物理建模或自适应强度控制来提升鲁棒性。总体而言，跨域视觉结果说明本文方法虽然在全参考指标上并非全面占优，但在实际主观观感和清晰度改善方面仍具有较强应用价值。

\subsection{结果讨论}
综合同域与跨域实验可得出三点结论。第一，在同域 matched 条件下，MP-CycleGAN 在 PSNR/SSIM 上取得稳定的增益，说明灰世界先验、结构约束与感知约束能够在非配对训练下减少不必要的内容漂移与结构破坏。第二，在跨域测试中，全参考指标与无参考指标的“分歧”更明显：CycleGAN 更偏向参考一致性，而 MP-CycleGAN 更偏向感知质量提升与视觉稳健性，这提示水下增强在真实应用中应采用多指标联合评估，并根据下游任务需求选择更合适的优化目标。第三，消融实验揭示了各损失项的互补性：灰世界更偏颜色校正，结构一致性更偏结构保真，感知损失更偏语义与纹理稳定，颜色统计一致性更偏整体视觉质量提升但可能引入与参考不一致的风险。

为保证结论可信度，本文将评测脚本、关键输出文件与复现命令在附录中统一列出，使实验可复核；同时通过 bootstrap 置信区间降低小样本波动带来的偶然性。未来工作可进一步在更大规模的跨域数据上验证统计显著性，并结合下游任务（检测/分割）进行任务驱动评估，从而更全面地刻画水下增强方法的工程价值。

\subsection{效率与资源开销分析}
从工程部署角度，增强模型的推理开销直接影响其在水下机器人与在线视觉系统中的可用性。本文方法在推理阶段仅使用生成器 $G: A\to B$ 输出增强结果，其网络结构与标准 CycleGAN 的生成器一致，因此参数规模与推理复杂度与 CycleGAN 基线保持同量级。换言之，本文提出的多感知自监督约束主要作用于训练阶段，通过额外损失项引导生成器学习更稳定的映射；在推理阶段不需要额外的 VGG 特征提取或梯度约束计算，因而不会引入额外推理模块与显著延时。

在训练阶段，感知损失需要进行预训练网络前向传播，结构一致性损失需要计算梯度差分，均会增加一定训练时间与显存开销。该开销与训练批量大小、输入分辨率及 VGG 特征层深度相关。对于实际复现实验，建议在保证稳定收敛的前提下优先控制 \texttt{perceptual\_layer} 深度与损失权重，以在性能与训练成本之间取得平衡。

\subsection{误差来源与潜在失败模式}
水下增强任务的评估与解释通常受多种因素影响。首先，matched 子集的参考图像获取方式与成像风格会影响 PSNR/SSIM 的上限：当增强结果在亮度分布、色彩风格上更符合人类感知但与参考存在系统性差异时，PSNR/SSIM 可能降低。其次，水体类型、悬浮颗粒密度与光照分布变化会导致退化机理差异显著，非配对训练在跨域场景中更容易出现“增强强度不匹配”的现象。

结合本文的损失设计，可以预期以下情形更具挑战性：一是极端浑浊或强后向散射场景，局部对比度提升可能引入噪声放大与纹理伪影；二是强光照不均与大面积高亮区域，颜色校正约束可能导致局部过度白化或色彩漂移；三是含有细密重复纹理的区域，结构约束与对抗学习的博弈可能产生轻微振铃或边缘晕圈。为降低上述风险，可考虑引入自适应增强强度控制、对高光区域的约束项，或在训练中加入更丰富水体类型的数据以提升泛化稳定性。

\subsection{本章小结}
本章在统一评测脚本下对 MP-CycleGAN 进行了同域评测、跨域泛化测试与消融分析，并结合可视化结果解释了不同指标之间的差异来源。总体而言，MP-CycleGAN 在同域 matched 设置下对像素一致性与结构保真带来小幅但稳定的提升；在跨域设置下，其感知质量相关指标具有一定优势，但仍存在与参考一致性指标的权衡。上述结论为后续工作提供了明确方向：在保持推理开销不变的前提下，进一步提升跨域鲁棒性与增强强度自适应能力。

\section{结论与展望}
\subsection{结论}
本文围绕“真实水下场景难以获取严格配对数据、模型跨域泛化不稳定”这一核心问题，提出了物理与感知双重驱动的非配对水下图像增强方法 MP-CycleGAN。该方法以 CycleGAN 为基础框架，在保持非配对学习优势的同时，引入灰世界物理先验、结构梯度一致性与深层感知约束，旨在缓解无监督风格迁移常见的颜色漂移、结构破坏与伪影问题，并提升跨域稳定性。

在统一评测脚本下，本文对 EUVP 同域 matched 子集与 UIEB 跨域 matched 子集进行了系统评测，并通过 bootstrap 置信区间提升结果可复核性。实验结果显示，MP-CycleGAN 在同域 matched=200 的 PSNR/SSIM 上相较 CycleGAN 取得小幅提升（22.91 dB / 0.795 对比 22.82 dB / 0.783），同时在跨域 matched=890 条件下保持竞争性表现并在 UIQM 上略优（10.50 对比 10.48）。虽然其定量增益并不大，但结合消融与可视化结果可以看出，本文方法在颜色均衡、结构稳定性与视觉自然性方面具有更稳定的优势，这种优势更符合真实水下增强任务对鲁棒性与可用性的要求。消融实验进一步揭示了各损失项的作用方向与权衡关系，为后续方法改进提供了明确依据。

总体而言，本文工作在不增加推理复杂度的前提下，通过训练阶段的多感知自监督约束增强了非配对增强框架的可控性与可解释性，为无监督水下图像增强提供了一套可复现的工程实现与分析范式。

\subsection{展望}
后续可从以下方向继续改进：
\begin{enumerate}
    \item 引入更强的水下物理成像模型或可学习的颜色校正模块；
    \item 针对无参考指标与全参考指标的冲突，引入多目标优化或感知一致性评价；
    \item 在更大规模、更复杂水体类型的数据上进行验证，并结合下游任务（检测/分割）进行端到端评估；
    \item 增加对失败案例与不确定性因素的建模，例如对极端浑浊/强光照不均场景引入自适应增强强度控制；
    \item 在训练与推理效率方面，探索轻量化生成器结构或蒸馏策略，以满足资源受限设备的实时需求。
\end{enumerate}

%-----------------------------------------------------------
% 参考文献
%-----------------------------------------------------------
\newpage
% 参考文献：序号左顶格，数字加方括号
\begin{thebibliography}{99}
\addcontentsline{toc}{section}{参考文献}

\bibitem{1} C. Li, C. Guo, W. Ren, R. Cong, J. Hou, S. Kwong, and D. Tao, “An Underwater Image Enhancement Benchmark Dataset and Beyond,” IEEE Trans. Image Process., vol. 29, pp. 4376–4389, 2020.
\bibitem{2} M. J. Islam, Y. Xia, and J. Sattar, “Fast Underwater Image Enhancement for Improved Visual Perception,” IEEE Robot. Autom. Lett., vol. 5, no. 2, pp. 3227–3234, Apr. 2020.
\bibitem{3} C. Li, S. Anwar, J. Hou, R. Cong, C. Guo, and W. Ren, “Underwater Image Enhancement via Medium Transmission-Guided Multi-Color Space Embedding,” IEEE Trans. Image Process., vol. 30, pp. 4985–5000, 2021.
\bibitem{4} J.-Y. Zhu, T. Park, P. Isola, and A. A. Efros, “Unpaired Image-to-Image Translation Using Cycle-Consistent Adversarial Networks,” in Proc. IEEE Int. Conf. Comput. Vis. (ICCV), 2017, pp. 2242–2251.
\bibitem{5} K. He, J. Sun, and X. Tang, “Single Image Haze Removal Using Dark Channel Prior,” in Proc. IEEE Conf. Comput. Vis. Pattern Recognit. (CVPR), 2009, pp. 1956–1963.
\bibitem{6} P. L. J. Drews Jr., E. R. Nascimento, S. S. C. Botelho, and M. F. M. Campos, “Underwater Depth Estimation and Image Restoration Based on Single Images,” IEEE Comput. Graph. Appl., vol. 36, no. 2, pp. 24–35, 2016.
\bibitem{7} G. Buchsbaum, “A Spatial Processor Model for Object Colour Perception,” J. Franklin Inst., vol. 310, no. 1, pp. 1–26, 1980.
\bibitem{8} Z. Wang, A. C. Bovik, H. R. Sheikh, and E. P. Simoncelli, “Image Quality Assessment: From Error Visibility to Structural Similarity,” IEEE Trans. Image Process., vol. 13, no. 4, pp. 600–612, 2004.
\bibitem{9} J. Johnson, A. Alahi, and L. Fei-Fei, “Perceptual Losses for Real-Time Style Transfer and Super-Resolution,” in Proc. Eur. Conf. Comput. Vis. (ECCV), 2016.
\bibitem{10} K. Simonyan and A. Zisserman, “Very Deep Convolutional Networks for Large-Scale Image Recognition,” in Proc. Int. Conf. Learn. Represent. (ICLR), 2015.
\bibitem{11} M. Yang and A. Sowmya, “An Underwater Color Image Quality Evaluation Metric,” IEEE Trans. Image Process., vol. 24, no. 12, pp. 6062–6071, 2015.
\bibitem{12} K. Panetta, C. Gao, and S. Agaian, “Human-Visual-System-Inspired Underwater Image Quality Measures,” IEEE J. Oceanic Eng., vol. 41, no. 3, pp. 541–551, 2016.

\end{thebibliography}

%-----------------------------------------------------------
% 英文部分
%-----------------------------------------------------------
\newpage

% 英文题目: Times New Roman, 小四号
\begin{center}
    {\zihao{-4}\bfseries MP-CycleGAN: Multi-Perceptual CycleGAN for Unpaired Underwater Image Enhancement}
\end{center}

\vspace{\baselineskip}

% 英文摘要
\noindent\hspace{2em}{\zihao{-4}\bfseries Abstract:}
{\zihao{-4}\setstretch{1.0}
Underwater image enhancement is a critical preprocessing step for marine engineering and underwater vision tasks. Due to wavelength-dependent absorption and backscattering, underwater images often suffer from severe color casts, low contrast, and haze-like artifacts, which substantially degrade the performance of downstream perception algorithms. While supervised deep models can achieve strong quantitative results, they heavily rely on pixel-aligned paired data that is expensive or even infeasible to collect in real underwater environments, and they may generalize poorly across domains.

To address these limitations, this thesis proposes \textbf{MP-CycleGAN} (Multi-Perceptual CycleGAN), an unpaired underwater enhancement framework built upon CycleGAN. We incorporate three self-supervised constraints: a gray-world color prior to mitigate underwater color casts, a gradient-based structure consistency loss to reduce structural distortions, and a VGG-16 perceptual loss to preserve semantic content and fine textures. We evaluate the method under a unified inference and evaluation pipeline on EUVP and UIEB, and report bootstrap confidence intervals for key metrics to improve reproducibility.

Results show that MP-CycleGAN achieves competitive performance under both in-domain evaluation and cross-domain testing, with clear trends supported by ablation studies.
\par}

\vspace{\baselineskip}

% 英文关键词
\noindent\hspace{2em}{\zihao{-4}\bfseries Keywords:}
{\zihao{-4} underwater image enhancement; unpaired learning; CycleGAN; gray-world prior; perceptual loss; structure consistency}

%-----------------------------------------------------------
% 附录
%-----------------------------------------------------------
\newpage
\appendix
\section{复现说明与命令}
本附录给出关键结果的来源文件与可直接执行的命令模板，便于复核。输出的汇总表位于：\path{results/benchmarks/summary.csv}。

\subsection{指标来源（关键 summary.json）}
\begin{itemize}
    \item EUVP 同域（matched=200）
    \begin{itemize}
        \item CycleGAN：\newline\path{results/benchmarks/euvp/models/euvp_cyclegan_full/epoch_200/summary.json}
        \item MP-CycleGAN：\newline\path{results/benchmarks/euvp/models/euvp_mpcgan_stage2_s0/epoch_202/summary.json}
    \end{itemize}
    \item UIEB 跨域（matched=890）
    \begin{itemize}
        \item CycleGAN：\newline\path{results/benchmarks/uieb/models/euvp_cyclegan_full/epoch_200/summary.json}
        \item MP-CycleGAN：\newline\path{results/benchmarks/uieb/models/euvp_mpcgan_stage2_s0/epoch_202/summary.json}
    \end{itemize}
\end{itemize}

\subsection{推理与评测命令（Windows / PowerShell）}
1) EUVP 单向推理（A$\to$B）：
\begin{lstlisting}[language=bash]
python test.py `
  --dataroot "<EUVP_TEST_INP>" `
  --name "euvp_cyclegan_full" `
  --model test `
  --dataset_mode single `
  --model_suffix "_A" `
  --epoch 200 `
  --num_test 200
\end{lstlisting}

2) EUVP 评测：
\begin{lstlisting}[language=bash]
python scripts/evaluate_euvp_psnr_ssim.py `
  --inp_dir "<EUVP_TEST_INP>" `
  --gtr_dir "<EUVP_TEST_GTR>" `
  --pred_dir "results\\euvp_cyclegan_full\\test_200\\images" `
  --bootstrap_iters 2000 `
  --seed 123
\end{lstlisting}

\subsection{论文图表生成}
在仓库根目录运行：
\begin{lstlisting}[language=bash]
python scripts/plot_unified_paper_figures.py
\end{lstlisting}
图表输出目录：\texttt{docs/figures/}。

\end{document}
